\chapter{Famílias e igrejas resilientes}
\noindent\textit{Vinheta.} Um bairro enfrenta 72 horas de instabilidade em energia e pagamentos. Duas casas ficam em pânico; outra segue um plano simples: água, comida, lista de vizinhos vulneráveis, ponto de encontro e culto doméstico. A diferença não foi sorte: foi \textbf{preparo}.

\section{Mapa bíblico}
\begin{itemize}[leftmargin=*,noitemsep]
  \item \textbf{Mordomia} (Pv~6:6--8; 1~Tm~5:8): prover para os seus com diligência.
  \item \textbf{Comunhão} (At~2:42--47): partilha, oração e mesa comum.
  \item \textbf{Caridade ordenada} (Tg~1:27; Gl~6:10): preferência pelos vulneráveis.
  \item \textbf{Paz da cidade} (Jr~29:7): o bem do bairro é o nosso bem.
\end{itemize}

\begin{nicbox}
\textbf{NIC aplicado}\\
\textbf{Narrativa}: não vivemos de pânico, mas de esperança e serviço.\\
\textbf{Identidade}: somos família de Deus, não consumidores solitários.\\
\textbf{Controle}: quando portas fecham, abrimos portas de generosidade.
\end{nicbox}

\section{Checklists NIC --- Tempo de paz (família)}
\begin{checklist}[title={Paz: 7 hábitos básicos},breakable]
\begin{itemize}[leftmargin=*,noitemsep]
  \item \textbf{Narrativa} — \emph{Palavra primeiro}: 10 min/dia de leitura e oração (cap.~13).
  \item \textbf{Identidade} — Documentos organizados (físico e digital); contatos impressos.
  \item \textbf{Controle} — \textbf{Múltiplos meios de pagamento}: dinheiro trocado + Pix + cartão.
  \item \textbf{Água \& alimento} — 7 dias de itens simples (giram todo mês); gás/álcool/lenha \emph{seguros}.
  \item \textbf{Saúde} — remédios contínuos (30 dias), receitas impressas, kit primeiros socorros.
  \item \textbf{Energia \& luz} — lanternas/pilhas, \emph{power banks}, filtro de linha; gerador \emph{apenas em área ventilada}.
  \item \textbf{Comunicação} — plano de pontos de encontro; vizinhos vulneráveis mapeados.
\end{itemize}
\end{checklist}

\section{Checklists NIC --- Tempo de crise (família)}
\begin{checklist}[title={Crise: 48--72h},breakable]
\begin{itemize}[leftmargin=*,noitemsep]
  \item \textbf{Narrativa} — confirmar fatos antes de agir; 1 boletim por período (manhã/noite).
  \item \textbf{Identidade} — guardar documentos e dinheiro; avisar 2 contatos de referência.
  \item \textbf{Controle} — compras pequenas e fracionadas; recibos manuais; evitar deslocamentos longos.
  \item \textbf{Água \& alimento} — priorizar perecíveis; cozinhar em blocos; partilhar com vizinhos vulneráveis.
  \item \textbf{Saúde} — separar medicações do dia; registrar sinais de alerta (idosos/diabéticos/crianças).
  \item \textbf{Energia \& luz} — racionar bateria; desconectar eletrônicos sensíveis; segurança com fogo.
  \item \textbf{Comunicação} — ponto de encontro; recados escritos no condomínio/rua.
\end{itemize}
\end{checklist}

\section{Política de resiliência da igreja (tempo de paz)}
\begin{checklist}[title={5 frentes permanentes},breakable]
\begin{itemize}[leftmargin=*,noitemsep]
  \item \textbf{Formação} — trilha NIC anual (cap.~1--4) e oficinas (cap.~3, 6, 10, 11).
  \item \textbf{Rede de cuidado} — líderes de quarteirão; cadastro voluntário de habilidades.
  \item \textbf{Infra mínima} — água, extintores, kit primeiros socorros, impressora, lanternas, rádio simples.
  \item \textbf{Comunicação} — modelos de boletim; rota de decisões (quem fala, quando, onde).
  \item \textbf{Fundo solidário} — ajuda imediata (alimentos, remédios, transporte) com prestação de contas trimestral.
\end{itemize}
\end{checklist}

\section{Plano da igreja para 48--72h (tempo de crise)}
\begin{checklist}[title={Operação básica da comunidade},breakable]
\begin{itemize}[leftmargin=*,noitemsep]
  \item \textbf{Dia 1} — reconhecimento: quem está em risco; o que falta; qual o ponto de apoio.
  \item \textbf{Dia 2} — coordenação: escala de cozinha solidária, coleta de doações, compra coletiva em feiras/atacadistas.
  \item \textbf{Dia 3} — resiliência: visitas a enfermos/idosos; culto comunitário de encorajamento; registro de lições.
\end{itemize}
\end{checklist}

\section{Cenários especiais e adaptações}
\subsection*{Famílias com bebês/crianças}
Fraldas, leite/alternativos, remédios pediátricos, entretenimento sem tela, contatos de pediatra e vizinhos com carro.
\subsection*{Idosos e pessoas com deficiência}
Listas de medicação, andadores/cadeiras, rampas improvisadas, rede para checagem 2x/dia.
\subsection*{Interior/rural}
Poço/água de chuva com potabilização, estoques maiores, rádio e ferramentas manuais, rotas alternativas em estradas de terra.

\section{Perguntas para grupos pequenos}
\begin{itemize}[leftmargin=*,noitemsep]
  \item Qual item do \textbf{tempo de paz} você consegue implementar ainda esta semana?
  \item Que \textbf{papel} você pode assumir na rede de cuidado da igreja?
  \item Quem no seu quarteirão precisa de atenção especial em uma crise?
\end{itemize}

\bigskip
\noindent\textit{Próximo capítulo: Ética para quem constrói tecnologia — linhas vermelhas, objeção de consciência e design responsável.}
